\documentclass[11pt,a4paper]{article}
\usepackage[utf8]{inputenc}
\usepackage[T1]{fontenc}
\usepackage{amsmath,amssymb,amsthm,physics,geometry,booktabs,array,hyperref,mathtools}
\usepackage{xcolor}
\usepackage{colortbl}
\geometry{margin=2.5cm}
\hypersetup{colorlinks=true,linkcolor=black,citecolor=black,urlcolor=black}
\theoremstyle{plain}
\newtheorem{theorem}{Theorem}[section]
\newtheorem{corollary}[theorem]{Corollary}
\newtheorem{proposition}[theorem]{Proposition}
\theoremstyle{definition}
\newtheorem{definition}{Definition}[section]
\theoremstyle{remark}
\newtheorem{remark}{Remark}[section]
\newcommand{\lQ}{\ell_Q}
\newcommand{\sigt}{\sigma_t}
\newcommand{\dotv}{\dot\sigma_t}
\definecolor{lightgray}{gray}{0.93}
\definecolor{warnyellow}{RGB}{255,248,220}
\definecolor{warnborder}{RGB}{180,140,0}

\title{\textbf{Quantum Gravitational Dynamics}\\[4pt]
\Large Chapter 5: Cosmology}
\author{Romeo Matshaba}
\date{University of South Africa \\ \small DOI: 10.5281/zenodo.18827993}

\begin{document}
\maketitle

\begin{abstract}
In a homogeneous, isotropic universe the graviton field reduces to
$\sigma_\mu = (\sigma_t(t),0,0,0)$.  The master field equation
$\Box_g\sigma_\mu = Q_\mu + G_\mu + T_\mu + \kappa\lQ^2\Box_g^2\sigma_\mu$
becomes a fourth-order ODE in cosmic time, expanding the solution space
from the two-dimensional Friedmann phase space to four dimensions.  The
extra modes produce, in a unified framework: a candidate inflationary
mechanism at the Planck scale, a $w = -1$ dark energy equation of state
from kinetic energy alone, and an exact match to the observed dark energy
density $\rho_\Lambda \approx 6\times10^{-10}$~J/m$^3$ without a
cosmological constant.

\smallskip\noindent\textbf{Updated (March 2026).}  A companion
computation~\cite{matshaba2026:kin_var} performs the full variation of
the kinetic action $S_{\mathrm{kin}}$ tracking every $g(\sigma)$
dependence, identifying a potential term $V(\sigma_t)$ that may be
required for self-consistency of the Schwarzschild static sector.  The
consequences for cosmology are analysed in
Section~\ref{sec:kin_consequences}: the four-mode structure and
inflation candidate are unaffected; the exact $w=-1$ result becomes
conditional on whether the potential is required by the theory; and if
so, a specific $w(z)$ evolution emerges that is consistent with the
DESI~2024 mild deviation from $\Lambda$.
\end{abstract}

\tableofcontents
\newpage

%=============================================================
\section{Cosmological Symmetry Reduction}
\label{sec:symmetry}
%=============================================================

\subsection{The cosmological graviton field}

Homogeneity and isotropy uniquely fix the graviton field to have only a
temporal component:
\begin{equation}
  \sigma_\mu = (\sigma_t(t),\;0,\;0,\;0).
\end{equation}
Any nonzero spatial component would define a preferred direction, violating
isotropy.  The resulting metric is the flat FRW form:
\begin{equation}
  \dd s^2 = c^2\dd t^2 - a^2(t)\,\delta_{ij}\,\dd x^i\dd x^j.
\end{equation}

\subsection{Energy-momentum tensor of the cosmological $\sigma$-field}

The canonical stress tensor for this configuration gives:
\begin{equation}
  \rho_\sigma = \tfrac{1}{2}\dotv^2, \qquad
  p_\sigma = -\tfrac{1}{2}\dotv^2, \qquad
  \boxed{w_\sigma \equiv \frac{p_\sigma}{\rho_\sigma} = -1.}
  \label{eq:w}
\end{equation}
This $w=-1$ equation of state is mathematically exact for any nonzero
$\dotv$.  It arises purely from the kinetic structure of the
$\sigma$-field Lagrangian --- not from a potential, a cosmological
constant, or fine-tuning.

\begin{remark}[Conditionality of $w=-1$]
\label{rem:w_conditional}
The result~\eqref{eq:w} holds when the action contains only the kinetic
term $S_{\mathrm{kin}}$.  The full variation of $S_{\mathrm{kin}}$
(Section~\ref{sec:kin_consequences}) indicates a potential
$V(\sigma_t)$ may be required for self-consistency of the static sector.
If so, $w$ deviates from $-1$ by $\Delta w \approx 2V(\sigma_t)/\dotv^2$,
which is small at high redshift ($\sigma_t\ll 1$) and becomes
observationally relevant at low redshift ($\sigma_t \sim 1$).  See
Section~\ref{sec:w_modified} for the full analysis.
\end{remark}

%=============================================================
\section{The Fourth-Order Cosmological Equation}
\label{sec:fourth_order}
%=============================================================

The master field equation in the FRW background reduces to:
\begin{equation}
  \boxed{\ddot\sigma_t + 3H\dotv
  - \lQ^2\!\left[\ddddot\sigma_t + 3H\dddot\sigma_t
                 + 3\dot H\ddot\sigma_t
                 + (3H^2 + 3\dot H)\dotv\right] = S_t}
  \label{eq:cosmo_master}
\end{equation}
where $H = \dot a/a$, $\lQ = \ell_{\mathrm{Pl}} \approx 1.616\times10^{-35}$~m, and
\begin{equation}
  S_t = \frac{8\pi G}{c^4}Q_t[\sigma_t,\dotv]
      + \frac{4\pi G}{c^2}\rho_m\sigma_t, \qquad
  Q_t = \sigma_t\dotv^2 \quad\text{(homogeneous limit).}
\end{equation}
This is \textbf{fourth-order in time}.  The full state requires
$\{\sigma_t,\dotv,\ddot\sigma_t,\dddot\sigma_t\}$ plus $\{a,\dot a\}$,
a six-dimensional phase space reduced to five by the Friedmann
constraint.  Standard $\Lambda$CDM has only two dimensions.

\begin{remark}[Completeness of the FRW field equation]
The $Q_t = \sigma_t\dotv^2$ source is the homogeneous-limit analogue of
the static Schwarzschild $Q_t = \sigma_t(\partial_r\sigma_t)^2$.  The
companion calculation~\cite{matshaba2026:kin_var} shows the static
sector has additional terms from the geometric coupling $G_\mu$.  In
FRW, $G_\mu$ is encoded in the fourth-order stiffness bracket --- these
are physically distinct sectors.  However, a parallel full variation of
$S_{\mathrm{kin}}$ in the FRW background, tracking $\partial
g^{\mu\nu}/\partial\sigma_t$ and $\partial\sqrt{-g}/\partial\sigma_t$,
has not yet been performed.  The $Q_t$ form above is therefore the
leading-order expression; correction terms may arise at higher order
in $\sigma_t$.
\end{remark}

%=============================================================
\section{The Four-Mode Solution}
\label{sec:modes}
%=============================================================

\subsection{Vacuum de Sitter background}

For $H = H_0 = \mathrm{const}$ and $S_t = 0$, the equation factorises:
\begin{equation}
  (1 - \lQ^2\partial_t^2)(\partial_t^2 + 3H_0\partial_t)\sigma_t = 0.
\end{equation}
The general solution has four modes:
\begin{equation}
  \boxed{\sigma_t(t) = C_1 + C_2\,e^{-3H_0 t}
                     + C_3\,e^{+t/\lQ} + C_4\,e^{-t/\lQ}}
  \label{eq:four_modes}
\end{equation}

\begin{table}[h]
  \centering
  \caption{Mode structure and physical interpretation.
  The four-mode structure is a direct consequence of the
  quantum stiffness $\kappa\lQ^2\Box_g^2\sigma_\mu$ in the master
  equation and is independent of the $Q_t$ source details.}
  \label{tab:modes}
  \begin{tabular}{@{}clll@{}}
    \toprule
    Mode & Timescale & Character & Stability \\
    \midrule
    $C_1$ & $\infty$ & Constant background & Neutral \\
    $C_2\,e^{-3H_0t}$ & $\tau_H = 1/(3H_0) \sim 10^{17}$~s & Hubble decay & Stable \\
    $C_3\,e^{+t/\lQ}$ & $\lQ/c \sim 10^{-43}$~s & Exponential growth & Unstable \\
    $C_4\,e^{-t/\lQ}$ & $\lQ/c \sim 10^{-43}$~s & Rapid damping & Stable \\
    \bottomrule
  \end{tabular}
\end{table}

\subsection{The growing mode and nonlinear saturation}

The $C_3$ mode is the Ostrogradsky instability inherent to fourth-order
theories.  In the linear regime, after $t = 100\,\lQ/c \approx
5\times10^{-42}$~s:
$\sigma_t(t)/\sigma_t(0) \approx e^{100} \approx 2.7\times10^{43}$.

The nonlinear self-interaction $Q_t = \sigma_t\dotv^2$ saturates the
growth: the back-reaction (scaling as $e^{3t/\lQ}$) overtakes linear
growth ($e^{t/\lQ}$), giving:
\begin{equation}
  \sigma_{\mathrm{sat}} \sim \ell_{\mathrm{Pl}}.
\end{equation}
\emph{Status:} saturation is a conjecture requiring numerical integration.
The four-mode structure itself is exact and unaffected by the kinetic
variation results of Section~\ref{sec:kin_consequences}.

%=============================================================
\section{Dark Energy as Attractor}
\label{sec:dark_energy}
%=============================================================

\subsection{The constant-velocity solution}

With $\dotv = v = \mathrm{const}$, all higher derivatives vanish.
Self-consistency with the Friedmann equation gives:
\begin{equation}
  \boxed{v = \sqrt{\frac{3c^2}{4\pi G}}\,H_0, \qquad
         \rho_\sigma = \frac{3H_0^2 c^4}{8\pi G}, \qquad w = -1.}
  \label{eq:attractor}
\end{equation}

\subsection{The one-parameter de Sitter family}

A one-parameter family of exact solutions exists, parametrised by $v$:
\begin{align}
  \dotv &= v, \\
  H &= \sqrt{\frac{4\pi G}{3c^2}}\,v, \\
  a(t) &= a_0\,e^{Ht}.
\end{align}

\subsection{Exact match to observed dark energy}

For $H = H_0 \approx 2.2\times10^{-18}$~s$^{-1}$:
\begin{equation}
  \rho_\sigma = \tfrac{1}{2}\dotv^2 = \frac{3H_0^2 c^4}{8\pi G}
  \approx 5.3\times10^{-10}\;\mathrm{J/m^3}.
\end{equation}
The observed value is $\rho_\Lambda \approx 6\times10^{-10}$~J/m$^3$
(Planck~2018 \cite{Planck2018}).  This is a match within observational
uncertainty without any cosmological constant.

\begin{remark}[Attractor shift from potential]
\label{rem:attractor_shift}
If the potential $V(\sigma_t)$ identified in
Section~\ref{sec:kin_consequences} is present, the field equation
gains the term $-V'(\sigma_t) = -\sigma_t(1-\sigma_t^2)/4$, modifying
the balance condition.  The constant-velocity attractor~\eqref{eq:attractor}
still exists but at a shifted velocity:
\begin{equation}
  3H_0 v + V'(\sigma_t) = Q_t[\sigma_t, v] + \text{matter terms}.
\end{equation}
The correction to $\rho_\sigma$ is of order $V(\sigma_t)/\dotv^2$,
small when $\sigma_t \ll 1$.  A precise value requires numerical
integration of the full attractor equation.
\end{remark}

%=============================================================
\section{Modified Friedmann Equations}
\label{sec:friedmann}
%=============================================================

\begin{equation}
  \boxed{\begin{cases}
    H^2 = \dfrac{8\pi G}{3c^2}\!\left(\rho_m + \rho_r
          + \dfrac{1}{2}\dotv^2\right) \\[10pt]
    \dfrac{\ddot a}{a} = -\dfrac{4\pi G}{3c^2}\!\left(\rho_m + \rho_r
    + 3p_m + 3p_r - \dotv^2\right)
  \end{cases}}
  \label{eq:friedmann}
\end{equation}
Accelerated expansion ($\ddot a > 0$) occurs whenever
$\dotv^2 > \rho_m + \rho_r + 3p_m + 3p_r$.

%=============================================================
\section{Stability Analysis}
\label{sec:stability}
%=============================================================

\subsection{Saddle-point character}

Linearising around the constant-velocity attractor:
\begin{equation}
  r^2 + 3H_0 r - 6H_0^2 = 0 \quad\Longrightarrow\quad
  r_\pm = H_0\frac{-3 \pm \sqrt{33}}{2}.
\end{equation}
Numerically: $r_+ \approx +1.37\,H_0$ (unstable), $r_- \approx
-4.37\,H_0$ (stable).  The de Sitter solution is a saddle point.

\subsection{Perturbation decay}

Perturbations $\delta\sigma_k(t)$ around the attractor decay as:
\begin{equation}
  \delta\sigma_k(t) = A_k\,e^{-2Ht} + B_k\,e^{-3Ht} + C_k\,e^{-t/\lQ}.
\end{equation}
No growing perturbation mode.  The dark energy component is perfectly
smooth.

%=============================================================
\section{Three-Phase Cosmological Evolution}
\label{sec:phases}
%=============================================================

\begin{table}[h]
  \centering
  \caption{Three-phase QGD cosmological evolution.  The phase
  structure is robust to the potential correction identified in
  Section~\ref{sec:kin_consequences}, since $V(\sigma_t) \to 0$
  when $\sigma_t \to 0$ at early times.}
  \label{tab:phases}
  \begin{tabular}{@{}p{1.9cm}p{3.3cm}p{4.6cm}p{3.5cm}@{}}
    \toprule
    Phase & Epoch & Dominant mode & Outcome \\
    \midrule
    Quantum & $t \lesssim 10\,t_{\mathrm{Pl}}$ &
      $C_3 e^{+t/\lQ}$: growth, saturates near $\sigma_t \sim \ell_{\mathrm{Pl}}$ &
      Seeds inflation \\[4pt]
    Inflation & $10\,t_{\mathrm{Pl}} \lesssim t \lesssim t_{\mathrm{end}}$ &
      $\dotv \approx \mathrm{const}$; de Sitter &
      $N \sim 60$ e-folds \\[4pt]
    Classical & $t > t_{\mathrm{end}}$ &
      $C_1 + C_2 e^{-3Ht}$ &
      Late-time dark energy \\
    \bottomrule
  \end{tabular}
\end{table}

The energy exchange rate between $\sigma$-field and matter:
\begin{equation}
  \frac{\dd\rho_\sigma}{\dd t}
  = \frac{4\pi G}{c^2}\rho_m\sigma_t\dotv
  + \frac{8\pi G}{c^4}\sigma_t\dotv^3.
\end{equation}
At late times ($\rho_m \to 0$), self-interaction sustains constant
$\rho_\sigma$ at the attractor.

%=============================================================
\section{Consequences of the Kinetic Variation Analysis}
\label{sec:kin_consequences}
%=============================================================

The companion paper~\cite{matshaba2026:kin_var} performs the complete
symbolic variation of $S_{\mathrm{kin}}$ with respect to $\sigma_t$,
tracking the three contributions: direct variation, inverse-metric
variation $\partial g^{\mu\nu}/\partial\sigma_t$, and Jacobian variation
$\partial\sqrt{-g}/\partial\sigma_t$.  The key result for cosmology is
summarised here.

\subsection{The potential term}
\label{sec:potential}

The full Euler--Lagrange equation from $S_{\mathrm{kin}}$ in the
Schwarzschild static sector requires a source
\begin{equation}
  \mathcal{S} = \frac{\sigma_t(\sigma_t^2 - 1)}{4r^2}
  = -\frac{\partial V}{\partial\sigma_t}\cdot\frac{1}{r^2},
\end{equation}
where the potential is
\begin{equation}
  \boxed{V(\sigma_t) = \frac{\sigma_t^2}{8} - \frac{\sigma_t^4}{16}
  = \frac{\sigma_t^2(2-\sigma_t^2)}{16}.}
  \label{eq:potential}
\end{equation}
This potential vanishes at $\sigma_t = 0$ (flat space) and is
maximised between $\sigma_t = 0$ and $\sigma_t = 1$ (horizon).  It
connects to the NJL chiral symmetry-breaking structure of the
microscopic Dirac theory~\cite{matshaba2026:kin}.

There are two paths to resolving the static sector:

\begin{description}
  \item[Path 1 (geometric):] The geometric coupling $G_t$ from the
    Kerr--Schild sector supplies the missing source entirely.
    No potential term is needed.  In this case $V = 0$ in cosmology.

  \item[Path 2 (potential):] A potential term $S_{\mathrm{pot}} =
    -\int\sqrt{-g}\,V(\sigma_t)\,d^4x$ is present in the action.
    In this case $V \neq 0$ in cosmology.
\end{description}

Which path is correct is determined by an explicit symbolic computation
of $G_t$ from the Kerr--Schild structure (currently open).

\subsection{Consequences for $w = -1$}
\label{sec:w_modified}

\paragraph{Path 1 ($V = 0$):}
The canonical stress tensor gives $\rho_\sigma = \tfrac{1}{2}\dotv^2$,
$p_\sigma = -\tfrac{1}{2}\dotv^2$, and $w = -1$ exactly, as derived in
Section~\ref{sec:symmetry}.  Chapter~5 stands without modification.

\paragraph{Path 2 ($V \neq 0$):}
For the QGD phantom-type kinetic structure with $L \propto
+\tfrac{1}{2}\dotv^2$, adding $V$ gives
\begin{equation}
  \rho_\sigma = \tfrac{1}{2}\dotv^2 - V(\sigma_t), \qquad
  p_\sigma    = -\tfrac{1}{2}\dotv^2 - V(\sigma_t),
\end{equation}
so that
\begin{equation}
  \boxed{w = \frac{p_\sigma}{\rho_\sigma}
  = -\frac{\tfrac{1}{2}\dotv^2 + V(\sigma_t)}
          {\tfrac{1}{2}\dotv^2 - V(\sigma_t)}.}
  \label{eq:w_modified}
\end{equation}
For $V \ll \tfrac{1}{2}\dotv^2$, this gives
\begin{equation}
  w \approx -1 - \frac{4V(\sigma_t)}{\dotv^2}.
  \label{eq:w_approx}
\end{equation}

\begin{proposition}[Redshift evolution of $w$]
\label{prop:wz}
If the potential~\eqref{eq:potential} is present, $w(z)$ evolves with
redshift through $\sigma_t(z)$:
\begin{equation}
  w(z) \approx -1 - \frac{4V(\sigma_t(z))}{\dotv^2(z)}.
  \label{eq:wz}
\end{equation}
Two limiting regimes:
\begin{enumerate}
  \item[\rm(i)] High redshift ($\sigma_t \ll 1$, early universe):
    $V(\sigma_t) \approx \sigma_t^2/8 \to 0$, so $w(z) \to -1$.
    The dark energy equation of state is indistinguishable from
    a cosmological constant at high redshift.
  \item[\rm(ii)] Low redshift ($\sigma_t \sim 1$, today):
    $V \sim 1/16$ and the deviation from $-1$ is observable.
    For the attractor solution~\eqref{eq:attractor},
    $\dotv^2 = 3c^2 H_0^2/(4\pi G)$, giving:
    \begin{equation}
      \Delta w \equiv w + 1 \approx
      -\frac{4V(\sigma_t^{(0)})}{3c^2H_0^2/(4\pi G)}.
      \label{eq:deltaw}
    \end{equation}
    The magnitude of $\Delta w$ depends on $\sigma_t^{(0)}$ (the
    present value of $\sigma_t$), which must be determined from
    initial conditions.
\end{enumerate}
\end{proposition}

\subsection{Relation to DESI 2024}
\label{sec:desi}

The DESI collaboration~\cite{DESI2024} reports mild evidence for a
dynamical dark energy equation of state:
\begin{equation}
  w_0 \approx -0.95, \qquad w_a \approx -0.75 \quad
  \text{(2--3$\sigma$ deviation from } w=-1\text{)}.
\end{equation}
A pure cosmological constant predicts $w = -1$ exactly, which is in
tension with this result at the $2$--$3\sigma$ level.

The QGD prediction under Path~2 is qualitatively consistent with the
DESI signal:
\begin{enumerate}
  \item $w(z) \to -1$ at high redshift (consistent with CMB and BAO
    at $z \gtrsim 0.5$),
  \item $w(z)$ deviates from $-1$ at low redshift through the
    growth of $V(\sigma_t(z))$.
\end{enumerate}
This is precisely the shape of the DESI preferred region in the
$w_0$--$w_a$ plane.

\begin{remark}[Falsifiability]
Path~1 predicts $w = -1$ exactly and is falsified if the DESI
deviation from $\Lambda$ is confirmed at $> 3\sigma$.
Path~2 predicts a specific $w(z)$ from the single function $V(\sigma_t)
= \sigma_t^2/8 - \sigma_t^4/16$ with no free parameters beyond the
initial condition $\sigma_t^{(0)}$.  Once $\sigma_t^{(0)}$ is fixed
by matching to $\rho_\Lambda^{\mathrm{obs}}$, the entire $w(z)$ history
is a zero-parameter prediction.  This is testable by DESI Year~5
and Euclid.
\end{remark}

\subsection{Impact on attractor and dark energy density}

For Path~2, the constant-velocity balance condition~\eqref{eq:attractor}
is modified.  With $V'(\sigma_t) = \sigma_t(1-\sigma_t^2)/4$, the
attractor equation becomes:
\begin{equation}
  3H_0 v + \frac{\sigma_t(1-\sigma_t^2)}{4} = Q_t[\sigma_t, v]
  + \text{matter/radiation terms.}
  \label{eq:attractor_modified}
\end{equation}
This shifts the attractor velocity $v$ by an amount of order
$V'(\sigma_t)/(3H_0)$.  The implied correction to the dark energy
density is:
\begin{equation}
  \delta\rho_\sigma \sim V(\sigma_t) \lesssim \frac{1}{16}
  \quad [\text{in natural units where } \dotv = 1].
\end{equation}
Numerically, for $\sigma_t^{(0)} \sim 0.1$--$0.3$ at the present
epoch, $V(\sigma_t^{(0)}) \ll \tfrac{1}{2}\dotv^2$ and the dark energy
density match of Section~\ref{sec:dark_energy} is preserved to within
a few percent.

\subsection{Summary of consequences}

\begin{table}[h]
  \centering
  \caption{Impact of kinetic variation results on the cosmology chapter.
  Path~1: geometric coupling $G_t$ closes the Schwarzschild gap.
  Path~2: potential $V(\sigma_t)$ is required.}
  \label{tab:consequences}
  \begin{tabular}{@{}p{4.8cm} p{3.6cm} p{3.6cm}@{}}
    \toprule
    Result & Path 1 ($G_t$ closes gap) & Path 2 ($V$ needed) \\
    \midrule
    \rowcolor{lightgray}
    Four-mode solution (Eq.~\ref{eq:four_modes})
      & Unchanged & Unchanged \\
    Inflation candidate ($C_3$ mode)
      & Unchanged & Unchanged \\
    \rowcolor{lightgray}
    Perturbation decay
      & Unchanged & Unchanged \\
    $w = -1$ exactly
      & Yes (exact) & No: $w(z)$ from Eq.~\eqref{eq:wz} \\
    \rowcolor{lightgray}
    $\rho_\Lambda$ matching
      & Exact (Eq.~\ref{eq:attractor}) & Slight shift (Eq.~\ref{eq:attractor_modified}) \\
    DESI $w_0$--$w_a$ tension
      & Present ($2$--$3\sigma$) & Naturally resolved \\
    \rowcolor{lightgray}
    New prediction
      & None & $w(z)$ from $V(\sigma_t)$ \\
    \bottomrule
  \end{tabular}
\end{table}

%=============================================================
\section{Testable Predictions}
\label{sec:predictions}
%=============================================================

\subsection{CMB acoustic peak shift}

If $\sigma_t$ contributes a fraction $f_\sigma$ of the critical density
at recombination:
\begin{equation}
  \frac{\Delta\ell_1}{\ell_1} \approx -\frac{f_\sigma}{2}.
\end{equation}
Current precision: $\ell_1 = 220.8 \pm 0.4$ ($\sim 0.2\%$).  Detectable
if $f_\sigma > 0.004$.

\subsection{Dark energy equation of state}

\begin{itemize}
  \item \textbf{Path~1:} $w = -1$ exactly.  Any confirmed deviation at
    $> 3\sigma$ from DESI or Euclid \emph{falsifies} Path~1.
  \item \textbf{Path~2:} $w(z)$ follows Eq.~\eqref{eq:wz} with the
    universal potential $V(\sigma_t) = \sigma_t^2/8 - \sigma_t^4/16$.
    Once $\sigma_t^{(0)}$ is fixed, the full $w(z)$ is a
    zero-free-parameter prediction.  Current DESI preference for
    $w_0 \approx -0.95$ is consistent with $\sigma_t^{(0)} \sim 0.2$
    to order of magnitude.
\end{itemize}

\subsection{Inflation observables}

A full computation of the spectral index $n_s$ and tensor-to-scalar
ratio $r$ requires perturbation theory around the saturated $C_3$
background.  If the potential~\eqref{eq:potential} is present, it
contributes to the slow-roll parameters:
\begin{equation}
  \epsilon_V = \frac{c^2}{16\pi G}\left(\frac{V'}{V}\right)^2, \qquad
  \eta_V = \frac{c^2}{8\pi G}\frac{V''}{V}.
\end{equation}
For $V = \sigma_t^2/8 - \sigma_t^4/16$ near $\sigma_t \to 0$:
$V' \approx \sigma_t/4$, $V \approx \sigma_t^2/8$, giving $\epsilon_V
= c^2/(32\pi G)$ (independent of $\sigma_t$ at leading order).
The spectral tilt is then $n_s - 1 = -6\epsilon_V + 2\eta_V$.
A full computation is an open calculation.

%=============================================================
\section{Open Questions}
\label{sec:open}
%=============================================================

\paragraph{Established:}
\begin{itemize}
  \item Fourth-order ODE structure and four-mode solution
  \item $w = -1$ for kinetic-only action (exact)
  \item Attractor matches $\rho_\Lambda$ within observational error
  \item Linear perturbations around the attractor decay
\end{itemize}

\paragraph{Requiring calculation (updated):}
\begin{itemize}
  \item \textbf{[New]} Explicit computation of $G_t$ from the
    Kerr--Schild sector to determine whether Path~1 or Path~2 applies.
    This single calculation settles whether $w = -1$ exactly or
    whether $w(z)$ evolves.
  \item \textbf{[New]} Full variation of $S_{\mathrm{kin}}$ in the FRW
    background to verify the cosmological field equation~\eqref{eq:cosmo_master}
    from first principles.
  \item \textbf{[New]} Numerical computation of $\sigma_t^{(0)}$ from
    the modified attractor equation~\eqref{eq:attractor_modified} if
    Path~2.
  \item Nonlinear saturation amplitude (numerical ODE integration)
  \item Mechanism enforcing $C_3 \to 0$ if inflation from $C_3$ is used
  \item Spectral index and tensor-to-scalar ratio from slow-roll with $V$
  \item Boltzmann code implementation for CMB comparison
\end{itemize}

%=============================================================
\section{Conclusions}
\label{sec:conclusions}
%=============================================================

The fourth-order $\sigma$-field cosmology provides a unified framework
for inflation and dark energy from the same graviton field.  The
companion kinetic variation analysis~\cite{matshaba2026:kin_var}
identifies a structural choice that bifurcates the cosmological
predictions.

If the Kerr--Schild geometric coupling $G_t$ closes the Schwarzschild
self-consistency gap (Path~1), no potential term is needed, $w = -1$
exactly, and the cosmology of this chapter is unchanged.

If a potential term $V(\sigma_t) = \sigma_t^2/8 - \sigma_t^4/16$ is
required (Path~2), the consequences are:

\begin{enumerate}
  \item The four-mode structure and inflation candidate are unaffected.
  \item $w$ evolves with redshift via $w(z) \approx -1 -
    4V(\sigma_t(z))/\dotv^2(z)$, with $w \to -1$ at high $z$ and
    a deviation at low $z$.
  \item The dark energy density match $\rho_\sigma \approx \rho_\Lambda$
    is preserved to within a few percent for $\sigma_t^{(0)} \lesssim
    0.3$.
  \item The $w(z)$ prediction is consistent with the DESI~2024 mild
    deviation from $\Lambda$ and constitutes a zero-free-parameter
    prediction once $\sigma_t^{(0)}$ is fixed.
\end{enumerate}

The pivot calculation is explicit: compute $G_t$ from
$H_{\mathrm{KS}}(\ell\cdot\nabla)^2\sigma_t + (\nabla_\alpha\sigma_t)
\nabla^\alpha(H_{\mathrm{KS}}\ell^t\ell_t)$ for the Schwarzschild
Kerr--Schild structure and verify whether it equals the gap
$\Delta$ of~\cite{matshaba2026:kin_var}.

\begin{thebibliography}{9}

\bibitem{matshaba2026}
R.~Matshaba,
\textit{Universal Origins: Quantum Gravitational Dynamics},
University of South Africa (2026).
DOI: \href{https://doi.org/10.5281/zenodo.18827993}{10.5281/zenodo.18827993}.

\bibitem{matshaba2026:kin_var}
R.~Matshaba,
\textit{Kinetic Term Variation in QGD: Self-Consistency of the
Schwarzschild Solution},
companion paper (March 2026), same DOI.

\bibitem{matshaba2026:kin}
R.~Matshaba,
\textit{QGD Chapter~1: Derivation of the kinetic term from the
interacting Dirac theory},
companion paper, same DOI.

\bibitem{Planck2018}
Planck Collaboration,
\textit{Planck 2018 results VI: Cosmological parameters},
Astron.\ Astrophys.\ \textbf{641}, A6 (2020).

\bibitem{DESI2024}
DESI Collaboration,
\textit{DESI 2024 VI: Cosmological Constraints from the Measurements
of Baryon Acoustic Oscillations},
arXiv:2404.03002 (2024).

\end{thebibliography}

\end{document}
